\documentclass[11pt,a4paper]{article}
\usepackage[T1]{fontenc}
\usepackage[utf8]{inputenc}
\usepackage{lmodern}
\usepackage[a4paper,margin=1in]{geometry}
\usepackage{setspace}
\usepackage{hyperref}
\usepackage{enumitem}
\usepackage{xcolor}

\hypersetup{
    colorlinks=true,
    linkcolor=blue,
    urlcolor=blue,
    pdftitle={Code Structure Overview: Reflections.py},
    pdfauthor={Project Documentation}
}

\setstretch{1.08}
\setlist[itemize]{leftmargin=1.25em,itemsep=0.25em,topsep=0.25em}
\setlist[enumerate]{leftmargin=1.4em,itemsep=0.25em,topsep=0.25em}

\title{Code Structure Overview of \texttt{Reflections.py}}
\author{}
\date{}

\begin{document}
\maketitle

\section*{Purpose and Architectural Style}
\texttt{Reflections.py} is a single-module Streamlit application that combines mathematical transformation logic with interactive visualization. Conceptually, the file follows a pipeline:

\begin{center}
Input parsing $\rightarrow$ Reflection/state computation $\rightarrow$ Geometry generation $\rightarrow$ Plot assembly $\rightarrow$ Streamlit rendering.
\end{center}

The design is procedural but layered. Functions are grouped by responsibility (mathematics, plotting helpers, UI parsing, figure construction, orchestration). This gives the practical benefits of modular code while keeping deployment simple (one Python file).

\section*{High-Level Module Decomposition}
The module can be understood in seven functional layers.

\subsection*{1) Mathematical Domain Primitives}
The first group of functions defines the mathematical objects and transformations used throughout the app.

\begin{itemize}
    \item \textbf{Exceptional sequence generation}: \texttt{generate\_exceptional\_sequences(...)} builds eight classes of recursively defined barycentric sequences.
    \item \textbf{Coordinate conversion}: \texttt{bary\_to\_cart(...)} maps barycentric coordinates into 2D simplex coordinates.
    \item \textbf{Normalization}: \texttt{renormalize(...)} rescales triples to sum to 1.
    \item \textbf{Reflection operators}: \texttt{reflect\_1}, \texttt{reflect\_2}, \texttt{reflect\_3}, plus a dispatcher map \texttt{digit\_actions}.
\end{itemize}

These functions form the immutable ``math core'' used by both analytics and visualization.

\subsection*{2) Reflection Workflow Helpers}
A second layer supports sequence execution and pattern extraction.

\begin{itemize}
    \item \texttt{apply\_sequence(...)} applies a reflection word (string of digits) to one point.
    \item \texttt{matching\_reflection\_steps(...)} identifies positions where specific 2- and 3-symbol patterns appear.
    \item \texttt{plot\_reflected\_line\_steps(...)} applies partial sequences to line endpoints and draws transformed segments.
\end{itemize}

This layer bridges discrete symbolic data (reflection strings) with geometric traces.

\subsection*{3) Conic and Line Geometry}
The geometric rendering basis is prepared here.

\begin{itemize}
    \item \texttt{Q} encodes the quadratic form used for the conic.
    \item \texttt{compute\_conic\_cart\_points\_analytic()} computes conic points analytically in a fixed setup (constant sampling and tolerance), then clips points to the simplex.
    \item \texttt{base\_lines(...)} and \texttt{add\_line(...)} generate and render line families.
    \item \texttt{step\_size(...)} controls line-density scaling for larger indices.
\end{itemize}

A key structural decision is that conic computation is now intentionally non-configurable to preserve a fixed visual definition.

\subsection*{4) Configuration and Styling Data Models}
Two dataclasses establish explicit contracts between orchestration and rendering:

\begin{itemize}
    \item \texttt{SimplexFigureConfig}: feature toggles, line limits, label sizes, exceptional sequence parameters.
    \item \texttt{SimplexFigureStyle}: marker dictionaries and naming for plotted categories.
\end{itemize}

This reduces parameter sprawl in plotting code and improves readability of function calls.

\subsection*{5) UI Utility Layer}
A compact utility layer manages formatting and parsing:

\begin{itemize}
    \item \texttt{parse\_dimension\_vectors(...)} and \texttt{parse\_sequences(...)} sanitize user input.
    \item \texttt{format\_output\_html(...)} and \texttt{should\_render\_output\_line(...)} handle textual result presentation.
    \item \texttt{get\_visualization\_options(...)} and \texttt{get\_exceptional\_sequence\_options(...)} build Streamlit controls.
    \item \texttt{get\_default\_figure\_style(...)} centralizes style defaults.
\end{itemize}

\subsection*{6) Rendering Engine}
\texttt{build\_simplex\_figure(...)} is the central composition function and the most important structural hub.

It layers figure elements in a deterministic order:
\begin{enumerate}
    \item simplex boundary and fundamental triangle,
    \item conic curve,
    \item optional line families and reflected lines,
    \item exceptional points,
    \item correct/wrong nodes and labels,
    \item optional polygons,
    \item input vectors and simplex-corner annotations.
\end{enumerate}

This separation is architecturally strong: business logic computes records first, then this function handles visual assembly from those records.

\subsection*{7) Application Orchestration}
\texttt{main()} controls runtime flow:

\begin{enumerate}
    \item initialize Streamlit page,
    \item read user text inputs and toggles,
    \item build shared config/style objects,
    \item iterate through each reflection sequence,
    \item update quiver orientation via a 4-state machine,
    \item classify steps as \emph{correct} or \emph{wrong},
    \item accumulate geometry and output text,
    \item render either per-sequence figures or one combined figure.
\end{enumerate}

The module ends with the standard entry-point guard invoking \texttt{main()}.

\section*{Data Flow and State Logic}
The most distinctive structural element is the quiver-state machine in \texttt{main()}. It tracks a finite state in \{(1,1), (0,1), (1,0), (0,0)\}. For each reflection digit:

\begin{itemize}
    \item the current dimension vectors are reflected,
    \item the state transitions according to reflection-specific rules,
    \item the step is bucketed into \emph{correct} or \emph{wrong},
    \item corresponding geometry and formatted output lines are collected.
\end{itemize}

This means the program is not only plotting transformed points; it is simultaneously running a discrete classification process that drives coloring, legend semantics, and textual reporting.

\section*{Why the Current Structure Works Well}
\begin{itemize}
    \item \textbf{Single-file practicality}: easy execution and deployment for exploratory math visualization.
    \item \textbf{Layered responsibilities}: helper functions are mostly side-effect free; rendering and orchestration are concentrated.
    \item \textbf{Reusable rendering core}: one figure-builder supports both per-sequence and combined views.
    \item \textbf{Explicit contracts}: dataclasses make plotting configuration predictable.
    \item \textbf{Traceability}: each reflection step contributes to both visual traces and textual diagnostics.
\end{itemize}

\section*{Potential Future Refactor (Optional)}
If the project grows, the file could be split into modules without changing behavior:

\begin{itemize}
    \item \texttt{math\_core.py}: reflections, barycentric utilities, conic/line math,
    \item \texttt{plotting.py}: figure assembly and drawing helpers,
    \item \texttt{ui.py}: Streamlit controls and output formatting,
    \item \texttt{app.py}: top-level orchestration.
\end{itemize}

The current code already has near-module boundaries, so this would be an organizational extraction rather than a logic rewrite.

\section*{Summary}
\texttt{Reflections.py} is a mathematically driven visualization app with a clear internal architecture despite being a single file. Its structure is best understood as a staged transformation pipeline with three strong cores: symbolic reflections, geometric computation, and compositional plotting. The state-machine-based classification gives the application analytical depth beyond simple drawing, and the dataclass-backed rendering contract keeps complexity manageable.

\end{document}
